% !TEX program = pdflatex
% ================================================================
% KV-Cache Continuity and the Limits of Intermediate Representations
% in Depth-Truncated Transformers
% ================================================================
\documentclass{article}

% ---- NeurIPS 2026 style (or use ICLR/ICML as needed) ----
\usepackage[final]{neurips_2026}  % For submission; change to [preprint] if needed

% ---- Packages ----
\usepackage[utf8]{inputenc}
\usepackage[T1]{fontenc}
\usepackage{hyperref}
\usepackage{url}
\usepackage{booktabs}
\usepackage{amsfonts}
\usepackage{amsmath}
\usepackage{amssymb}
\usepackage{mathtools}
\usepackage{nicefrac}
\usepackage{microtype}
\usepackage{graphicx}
\usepackage{xcolor}
\usepackage{subcaption}
\usepackage{multirow}
\usepackage{array}
\usepackage{float}
\usepackage{algorithm}
\usepackage{algorithmic}
\usepackage{enumitem}

% ---- Custom commands ----
\newcommand{\kl}{\bar{D}_{\mathrm{KL}}}
\newcommand{\Eattn}{\mathcal{E}_{\text{attn}}}
\newcommand{\Ereadout}{\mathcal{E}_{\text{readout}}}
\newcommand{\dt}{d_t}
\newcommand{\Khat}{\hat{K}}
\newcommand{\Vhat}{\hat{V}}
\DeclareMathOperator{\LN}{LN}
\DeclareMathOperator{\RoPE}{RoPE}
\DeclareMathOperator{\RMSNorm}{RMSNorm}
\DeclareMathOperator{\Softmax}{Softmax}
\DeclareMathOperator*{\argmax}{arg\,max}

\title{KV-Cache Continuity and the Limits of Intermediate Representations\\in Depth-Truncated Transformers}

\author{
  Anonymous Authors \\
  Under Review
}

\begin{document}

\maketitle

% ================================================================
% ABSTRACT
% ================================================================
\begin{abstract}
Depth-adaptive inference strategies for autoregressive Transformers---including early exit, layer skipping, and self-speculative decoding---promise substantial compute savings by terminating token processing before the final layer. However, prior work evaluates only end-to-end generation quality, conflating two fundamentally distinct failure modes: \emph{attention memory poisoning} ($\Eattn$), where approximate Key--Value cache entries corrupt downstream multi-head attention, and \emph{unembedding manifold mismatch} ($\Ereadout$), where intermediate hidden states fail to align with the terminal readout layer. We introduce a \textbf{teacher-forced causal disentanglement protocol} that isolates these error sources as independently measurable quantities. Across extensive experiments on Qwen2.5-3B (36 layers, GQA) and Llama-3.2-1B (16 layers), we establish three results:
(1)~\textbf{Native KV State Projection}---populating skipped-layer cache slots using each layer's own $W_K$, $W_V$, and pre-layer normalization applied to the exit hidden state---preserves attention continuity with zero training, crossing the 1-bit mutual information boundary ($\kl \le \ln 2$) at exit depths $\ge 61\%$ across both architectures.
(2)~\textbf{Standalone hard early exit is structurally non-viable}: even a fully-parameterized ($r{=}2048$, 8.4M parameters) bottleneck MLP calibration adapter recovers baseline perplexity only at the cost of eliminating all compute savings.
(3)~A \textbf{mechanistic Key-stable / Value-evolving asymmetry} explains both results: Key cosine similarity to oracle states remains $\ge 0.75$ across all skipped layers, while Value similarity collapses below 0.30 within 2--3 layers of the exit point. We release our disentanglement harness as a reusable diagnostic tool for the depth-adaptive inference community.
\end{abstract}

% ================================================================
% 1. INTRODUCTION
% ================================================================
\section{Introduction}
\label{sec:intro}

Autoregressive large language models (LLMs) generate tokens sequentially, each requiring a full forward pass through all $L$ Transformer layers. Depth-adaptive inference methods---early exit~\cite{schwartz2020right,xin2020deebert,elbayad2020depth}, layer skipping~\cite{del2023skipdecode,raposo2024mixture}, and self-speculative decoding~\cite{elhoushi2024layerskip,zhang2024draft}---aim to reduce this cost by terminating computation at an intermediate layer $\dt < L$ for ``easy'' tokens. Yet a persistent, poorly understood problem undermines all such approaches: \emph{what exactly goes wrong when a token exits early?}

Prior work reports only end-to-end degradation metrics---aggregate perplexity increase, BLEU score drops, or downstream task accuracy losses---without distinguishing between two fundamentally different failure modes:

\begin{itemize}[nosep,leftmargin=*]
    \item \textbf{Attention Memory Poisoning} ($\Eattn$): When a token exits at layer $\dt$, its Key and Value cache entries for layers $\dt{+}1$ through $L$ are either absent or approximate. All subsequent tokens attend to these corrupted entries, potentially cascading errors through the autoregressive generation process.
    \item \textbf{Unembedding Manifold Mismatch} ($\Ereadout$): The hidden state $h_{\dt}$ at the exit layer resides on an intermediate representation manifold $\mathcal{M}_{\dt}$ that may be fundamentally misaligned with the terminal unembedding matrix $W_{\text{unembed}}$, which was trained to decode representations from $\mathcal{M}_L$.
\end{itemize}

These two failure modes have distinct computational and structural properties, yet no prior work has measured them independently. This paper introduces a \textbf{teacher-forced causal disentanglement protocol} that cleanly separates $\Eattn$ from $\Ereadout$ through controlled ablation conditions, and uses this protocol to establish three principal results.

\paragraph{Contributions.}
\begin{enumerate}[nosep,leftmargin=*]
    \item \textbf{Native KV State Projection:} A training-free mechanism that fills KV-cache holes by applying each skipped layer's own projection weights and pre-layer normalization to the exit hidden state (Section~\ref{sec:method}). This preserves attention continuity for exit depths $\ge 61\%$ of network depth, with zero base-model fine-tuning.
    
    \item \textbf{Causal disentanglement protocol:} A three-condition, teacher-forced ablation methodology that isolates attention memory corruption from readout mismatch as independently measurable error sources, enabling precise attribution of generation degradation (Section~\ref{sec:protocol}).
    
    \item \textbf{Definitive negative result on standalone early exit:} We demonstrate through a log-spaced adapter capacity sweep ($r \in \{64, \ldots, 2048\}$) that intermediate-layer readout requires $\ge 8.4$M learned parameters to match baseline quality at 50\% depth---eliminating any compute savings---establishing this as a structural limitation rather than a training artifact (Section~\ref{sec:readout}).
    
    \item \textbf{Mechanistic explanation via Key/Value asymmetry:} Layer-wise cosine similarity measurements reveal that Key representations stabilize early ($\cos(K^*, \Khat) \ge 0.75$ across all skipped layers) while Value representations decay rapidly ($\cos(V^*, \Vhat) < 0.30$ within 2--3 layers), explaining why attention routing survives truncation but semantic readout does not (Section~\ref{sec:mechanism}).
\end{enumerate}

% ================================================================
% 2. RELATED WORK
% ================================================================
\section{Related Work}
\label{sec:related}

\paragraph{Early exit and depth-adaptive inference.}
Early exit strategies attach auxiliary classifiers at intermediate layers to allow ``confident'' inputs to skip remaining computation~\cite{xin2020deebert,schwartz2020right,elbayad2020depth}. SkipDecode~\cite{del2023skipdecode} extends this to autoregressive generation with uniform-depth truncation. Mixture-of-Depths~\cite{raposo2024mixture} routes individual tokens to different layers. CALM~\cite{schuster2022confident} introduces confidence-based adaptive computation. However, none of these works disentangle the distinct failure modes of KV-cache corruption versus readout mismatch; all report only combined end-to-end quality metrics.

\paragraph{KV-cache management.}
Orthogonal work addresses KV-cache memory through token eviction~\cite{zhang2023h2o} and sliding window attention~\cite{xiao2024efficient}, which reduce the \emph{sequence length} dimension rather than the \emph{depth} dimension we address. Our work tackles the complementary problem of maintaining KV-cache \emph{continuity} across skipped layers during depth-adaptive inference.

\paragraph{Speculative decoding.}
Standard speculative decoding uses a separate smaller draft model~\cite{leviathan2023fast,chen2023accelerating}. EAGLE~\cite{li2024eagle} and Medusa~\cite{cai2024medusa} use dedicated high-capacity draft heads. LayerSkip~\cite{elhoushi2024layerskip} explores self-speculative decoding via early exit but requires full upper-layer recomputation during verification. Our Native Projection approach eliminates this recomputation by maintaining continuous KV-cache entries.

\paragraph{Intermediate representation analysis.}
The Tuned Lens~\cite{belrose2023eliciting} and Logit Lens~\cite{nostalgebraist2020logitlens} project intermediate representations to vocabulary space for interpretability. Our calibration adapter ablation (Section~\ref{sec:readout}) extends this line of work by systematically characterizing the capacity requirements for intermediate-to-terminal manifold alignment.

% ================================================================
% 3. NATIVE KV STATE PROJECTION
% ================================================================
\section{Native KV State Projection}
\label{sec:method}

Consider an $L$-layer autoregressive Transformer with hidden dimension $d_{\text{model}}$. When a token exits at depth $\dt < L$, its hidden state is $h_{\dt} \in \mathbb{R}^{d_{\text{model}}}$. The KV-cache slots for layers $l \in \{\dt{+}1, \ldots, L\}$ remain empty, creating ``holes'' that corrupt attention computation for all subsequent tokens.

\subsection{Projection Mechanism}

We propose filling these holes using each skipped layer's \emph{own} projection weights. For every skipped layer $l > \dt$:
\begin{align}
    \Khat_l &= W_K^{(l)}\, \LN_1^{(l)}(h_{\dt}), \label{eq:kproj} \\
    \Vhat_l &= W_V^{(l)}\, \LN_1^{(l)}(h_{\dt}), \label{eq:vproj}
\end{align}
where $W_K^{(l)}, W_V^{(l)}$ are the frozen Key and Value projection matrices of layer $l$, and $\LN_1^{(l)}$ is the pre-attention layer normalization. When rotary position embeddings are used, they are applied to $\Khat_l$ as in standard inference:
\begin{equation}
    \Khat_l^{\text{(cache)}} = \RoPE_p\!\left(\Khat_l\right).
\end{equation}

\paragraph{Computational cost.}
For a model with $n_{\text{kv}}$ KV heads and head dimension $d_{\text{head}}$, the per-token projection cost for $(L - \dt)$ skipped layers is:
\begin{equation}
    \text{FLOPs}_{\text{proj}} = 2 \times (L - \dt) \times d_{\text{model}} \times (n_{\text{kv}} \times d_{\text{head}}) + \mathcal{O}(\text{LN}).
\end{equation}
For Qwen2.5-3B ($n_{\text{kv}} = 2$, $d_{\text{head}} = 128$), this is $\sim 0.5\%$ of a full layer's compute---negligible compared to the MLP and attention costs of the skipped layers.

\paragraph{Key insight.}
The projection reuses the existing frozen weights of the model; no new parameters are introduced and no training is required. This is fundamentally different from learning a separate KV approximation, as the projection inherits the geometric structure that each layer's attention was trained to operate over.

% ================================================================
% 4. CAUSAL DISENTANGLEMENT PROTOCOL
% ================================================================
\section{Causal Disentanglement Protocol}
\label{sec:protocol}

To attribute generation degradation to specific failure modes, we design a teacher-forced, three-condition ablation protocol. All conditions use the same evaluation dataset (WikiText-2~\cite{merity2017pointer}), the same sequence length ($T = 64$ tokens), and strict teacher forcing to eliminate autoregressive error accumulation.

\subsection{Experimental Conditions}

Given a target exit depth $\dt$, we evaluate three conditions against the full-depth oracle baseline $\mathcal{M}^*$ (PPL$^* = 5.76$ for Qwen2.5-3B; PPL$^* = 8.69$ for Llama-3.2-1B):

\begin{enumerate}[nosep,leftmargin=*]
    \item \textbf{Attention-Only (Native Projection):} The KV cache for layers $l > \dt$ is populated via Equations~\eqref{eq:kproj}--\eqref{eq:vproj}, but the readout uses the \emph{full-depth} terminal state $h_L$. This isolates $\Eattn$: any quality degradation is attributable solely to KV-cache approximation.
    
    \item \textbf{Readout-Only (Oracle KV):} The KV cache contains pristine oracle entries $(K^*, V^*)$ from a full forward pass, but the prediction is decoded from the intermediate state $h_{\dt}$ (optionally through a calibration adapter). This isolates $\Ereadout$: any degradation is attributable solely to intermediate readout misalignment.
    
    \item \textbf{Attention-Only (Direct Copy):} A negative control where KV entries for skipped layers are naively copied from the exit layer: $\Khat_l = K_{\dt}$, $\Vhat_l = V_{\dt}$ for all $l > \dt$. This establishes a floor demonstrating that layer-specific projections are necessary.
\end{enumerate}

\subsection{Evaluation Metrics}

For each condition, we report:
\begin{itemize}[nosep,leftmargin=*]
    \item \textbf{Perplexity (PPL):} Standard language modeling perplexity under teacher forcing.
    \item \textbf{Mean KL divergence} $\kl = \frac{1}{T}\sum_{t=1}^T D_{\mathrm{KL}}(P^*_t \| \hat{P}_t)$: Token-averaged forward KL from oracle to approximate next-token distributions.
    \item \textbf{Top-1 agreement:} Fraction of positions where $\argmax \hat{P}_t = \argmax P^*_t$.
\end{itemize}

\paragraph{Information-theoretic acceptance criterion.}
We define the \textbf{1-bit mutual information boundary} as the operating threshold:
\begin{equation}
    \kl \le \ln(2) \approx 0.6931\;\text{nats}.
    \label{eq:1bit}
\end{equation}
This threshold has a natural information-theoretic interpretation: the approximate distribution $\hat{P}$ preserves at least as much mutual information with the ground-truth target as a single fair coin flip would destroy. It serves as a principled, pre-registered acceptance criterion for when KV-cache approximation is ``good enough.''

\begin{figure}[t]
    \centering
    \includegraphics[width=0.95\textwidth]{figures/fig1_disentanglement_schematic.pdf}
    \caption{\textbf{Causal disentanglement protocol.} A Transformer truncated at exit depth $\dt$ produces three experimental conditions: \textcolor{green!60!black}{\textbf{Attention-Only}} isolates KV-cache error ($\Eattn$) using Native Projection with full-depth readout; \textcolor{orange}{\textbf{Readout-Only}} isolates readout error ($\Ereadout$) using oracle KV with intermediate readout; \textcolor{red}{\textbf{Direct Copy}} provides a negative control establishing that layer-specific projection is necessary.}
    \label{fig:protocol}
\end{figure}

% ================================================================
% 5. KV CONTINUITY RESULTS
% ================================================================
\section{Results: KV-Cache Continuity Boundary}
\label{sec:kv_results}

\subsection{Primary Disentanglement on Qwen2.5-3B}

Table~\ref{tab:main} presents the complete causal disentanglement matrix for three representative exit depths on Qwen2.5-3B (36 layers, $d_{\text{model}} = 2048$, GQA with 16 query / 2 KV heads).

\begin{table}[t]
\centering
\caption{\textbf{Causal disentanglement matrix} on Qwen2.5-3B ($\text{PPL}^* = 5.76$). Teacher-forced evaluation on WikiText-2 ($T=64$). Native Projection achieves near-oracle attention quality at $\ge 67\%$ depth with zero training, while readout-only conditions fail at all depths even with trained adapters.}
\label{tab:main}
\small
\begin{tabular}{@{}llcccc@{}}
\toprule
\textbf{Exit Depth} & \textbf{Condition} & \textbf{PPL} & $\boldsymbol{\Delta}$\textbf{PPL} & $\boldsymbol{\kl}$ \textbf{(nats)} & \textbf{Top-1} \\
\midrule
\multicolumn{6}{l}{\textit{Full-Depth Baseline} ($\mathcal{M}^*$)} \\
$L = 36$ & Oracle & \textbf{5.76} & 0.00 & 0.0000 & 100.0\% \\
\midrule
\midrule
\multicolumn{6}{l}{\textit{Exit Depth $L = 18$ (50\% Relative Depth)}} \\
 & Readout-Only (Zero-Shot) & 715,406 & +715,401 & 11.7805 & 1.3\% \\
 & Readout-Only (3k-step MLP, $r{=}512$) & 359.39 & +353.63 & 4.3931 & 27.8\% \\
 & Attn-Only (Direct Copy) & 3,088 & +3,082 & 6.3036 & 15.3\% \\
 & \textbf{Attn-Only (Native Proj)} & \textbf{13.33} & \textbf{+7.57} & \textbf{0.9604} & \textbf{66.4\%} \\
\midrule
\multicolumn{6}{l}{\textit{Exit Depth $L = 24$ (67\% Relative Depth)}} \\
 & Readout-Only (Zero-Shot) & 200,684 & +200,678 & 10.7481 & 2.4\% \\
 & Attn-Only (Direct Copy) & 3,450 & +3,444 & 6.0820 & 22.2\% \\
 & \textbf{Attn-Only (Native Proj)} & \textbf{8.04} & \textbf{+2.28} & \textbf{0.4365} & \textbf{78.6\%} \\
\midrule
\multicolumn{6}{l}{\textit{Exit Depth $L = 30$ (83\% Relative Depth)}} \\
 & Readout-Only (Zero-Shot) & 30,369 & +30,363 & 7.7959 & 19.6\% \\
 & Attn-Only (Direct Copy) & 10,274 & +10,268 & 5.8089 & 29.9\% \\
 & \textbf{Attn-Only (Native Proj)} & \textbf{6.45} & \textbf{+0.68} & \textbf{0.1049} & \textbf{89.2\%} \\
\midrule
\multicolumn{6}{l}{\textit{Exit Depth $L = 34$ (94\% Relative Depth)}} \\
 & Readout-Only (Zero-Shot) & 70.91 & +65.15 & 2.4512 & 57.9\% \\
 & Attn-Only (Direct Copy) & 9.58 & +3.82 & 0.5946 & 75.7\% \\
 & \textbf{Attn-Only (Native Proj)} & \textbf{5.94} & \textbf{+0.18} & \textbf{0.0194} & \textbf{96.0\%} \\
\bottomrule
\end{tabular}
\end{table}

\paragraph{Key observations.}
Native Projection crosses the 1-bit boundary (Eq.~\ref{eq:1bit}) between $L{=}18$ (50\% depth, $\kl = 0.9604$) and $L{=}22$ (61\% depth, $\kl = 0.6151$), establishing an operating boundary at approximately 61\% of network depth. At $L{=}24$ (67\%), the KL divergence drops to $0.4365$ nats with 78.6\% Top-1 agreement; at $L{=}30$ (83\%), it reaches $0.1049$ nats with 89.2\% agreement---near-transparent to downstream attention.

In stark contrast, the readout-only condition fails catastrophically at all depths, with zero-shot perplexities ranging from $70.91$ ($L{=}34$) to $715{,}406$ ($L{=}18$), confirming that intermediate hidden states are fundamentally misaligned with the terminal unembedding matrix.

\subsection{Continuous Depth Sweep}

Figure~\ref{fig:kl_depth} presents the complete depth sweep across 11 exit depths ($L \in \{6, 10, 14, 18, 22, 24, 26, 28, 30, 32, 34\}$), revealing a smooth, monotonically decreasing relationship between exit depth and KL divergence for Native Projection. Direct Copy KL remains consistently several-fold to over an order of magnitude higher across all depths (from $6.6\times$ at $L{=}18$ to $30\times$ at $L{=}34$), confirming that layer-specific coordinate transformation---not mere state repetition---is responsible for the quality preservation.

\begin{figure}[t]
    \centering
    \includegraphics[width=0.7\textwidth]{figures/fig2_kl_vs_depth.pdf}
    \caption{\textbf{Information-theoretic continuity boundary.} Mean KL divergence between oracle and approximate next-token distributions as a function of relative exit depth on Qwen2.5-3B. Native Projection (blue) crosses below the $\ln(2)$ threshold (dashed) at $\sim$61\% depth; Direct Copy (red) remains above $\ln(2)$ until between 89--94\% depth (Table~\ref{tab:full_sweep}: KL $= 2.47$ at $L{=}32$, $0.59$ at $L{=}34$). The 11-point continuous sweep confirms the $\ge$61\% boundary is a smooth phase transition, not an artifact of sparse sampling.}
    \label{fig:kl_depth}
\end{figure}

\subsection{Cross-Architecture Validation on Llama-3.2-1B}

To verify that our findings are not specific to Qwen2.5's architecture, GQA configuration, or RoPE frequency scaling, we replicate the full disentanglement protocol on \texttt{meta-llama/Llama-3.2-1B} (16 layers, MHA, $d_{\text{model}} = 2048$). Table~\ref{tab:llama} presents results across 7 exit depths.

\begin{table}[t]
\centering
\caption{\textbf{Cross-architecture validation} on Llama-3.2-1B ($\text{PPL}^* = 8.69$, 16 layers). Native Projection crosses the $\ln(2)$ boundary at 50\% depth ($L{=}8$), confirming architectural invariance. All values from teacher-forced evaluation on WikiText-2 ($T=64$).}
\label{tab:llama}
\small
\begin{tabular}{@{}cccccc@{}}
\toprule
\textbf{Exit Depth} & \textbf{Rel.\ Depth} & \textbf{Condition} & \textbf{PPL} & $\boldsymbol{\kl}$ \textbf{(nats)} & \textbf{Top-1} \\
\midrule
$L = 8$ & 50.0\% & Readout-Only & 147,615 & 9.8892 & 2.1\% \\
 & & Direct Copy & 1,026 & 4.6518 & 23.5\% \\
 & & \textbf{Native Proj} & \textbf{17.04} & \textbf{0.6081} & \textbf{69.6\%} \\
\midrule
$L = 10$ & 62.5\% & Direct Copy & 172.37 & 2.8304 & 37.6\% \\
 & & \textbf{Native Proj} & \textbf{14.83} & \textbf{0.4344} & \textbf{77.2\%} \\
\midrule
$L = 11$ & 68.8\% & Direct Copy & 157.37 & 2.5396 & 42.1\% \\
 & & \textbf{Native Proj} & \textbf{12.48} & \textbf{0.3414} & \textbf{81.7\%} \\
\midrule
$L = 12$ & 75.0\% & Direct Copy & 44.44 & 1.4249 & 52.6\% \\
 & & \textbf{Native Proj} & \textbf{11.79} & \textbf{0.2715} & \textbf{82.3\%} \\
\midrule
$L = 13$ & 81.3\% & Direct Copy & 27.42 & 0.9525 & 65.1\% \\
 & & \textbf{Native Proj} & \textbf{11.51} & \textbf{0.2360} & \textbf{83.1\%} \\
\midrule
$L = 14$ & 87.5\% & Direct Copy & 17.00 & 0.5959 & 70.6\% \\
 & & \textbf{Native Proj} & \textbf{10.58} & \textbf{0.1392} & \textbf{87.0\%} \\
\bottomrule
\end{tabular}
\end{table}

\paragraph{Architectural invariance.}
Remarkably, on Llama-3.2-1B, Native Projection already crosses the $\ln(2)$ boundary at 50\% depth ($L{=}8$, $\kl = 0.6081$), compared to $\sim$61\% on Qwen2.5-3B. Figure~\ref{fig:cross_arch} visualizes this cross-architecture comparison. This suggests that the operating boundary may be \emph{conservative} for some architectures, and that the structural property enabling KV continuity---the geometric stability of Key representations across depth---is a general property of pre-trained Transformers, not an artifact of any specific architecture.

\begin{figure}[t]
    \centering
    \includegraphics[width=0.95\textwidth]{figures/fig5_cross_architecture.pdf}
    \caption{\textbf{Cross-architecture validation.} (a)~KL divergence and (b)~Top-1 agreement as a function of relative exit depth for Native Projection on Qwen2.5-3B (36 layers) and Llama-3.2-1B (16 layers). Both architectures exhibit smooth, monotonically improving quality with depth. Llama-3.2-1B crosses the $\ln(2)$ boundary at a lower relative depth ($\sim$50\%) than Qwen2.5-3B ($\sim$61\%), confirming architectural invariance.}
    \label{fig:cross_arch}
\end{figure}

% ================================================================
% 6. STANDALONE READOUT IS NON-VIABLE
% ================================================================
\section{Results: Standalone Early Exit is Structurally Non-Viable}
\label{sec:readout}

\subsection{Calibration Adapter Design}

Following the Tuned Lens~\cite{belrose2023eliciting} paradigm, we train exit calibration adapters to align intermediate representations with the terminal manifold. Each adapter is a 2-layer bottleneck MLP with GELU activation and zero-initialized output projection (residual connection):
\begin{equation}
    h_{\text{calib}} = h_{\dt} + W_2 \cdot \text{GELU}\!\left(W_1 \, \LN(h_{\dt})\right),
    \label{eq:adapter}
\end{equation}
where $W_1 \in \mathbb{R}^{d_{\text{model}} \times r}$ and $W_2 \in \mathbb{R}^{r \times d_{\text{model}}}$, with bottleneck rank $r$. Training uses 3{,}000 steps of knowledge distillation against the full-depth model's output distribution, with the base model frozen.

\subsection{Capacity Scaling Ablation}

Table~\ref{tab:adapter} presents the adapter capacity sweep at $L{=}18$ (50\% depth), ranging from $r{=}64$ (268K parameters) to $r{=}2048$ (8.4M parameters---a full hidden-dimension bottleneck).

\begin{table}[t]
\centering
\caption{\textbf{Adapter capacity saturation} at $L{=}18$ on Qwen2.5-3B. Even at $r{=}2048$ (8.4M parameters), the adapter barely matches baseline PPL while Top-1 agreement plateaus at 81.7\%---far below the 100\% achievable by Native Projection with zero learned parameters.}
\label{tab:adapter}
\small
\begin{tabular}{@{}cccccl@{}}
\toprule
\textbf{Rank} ($r$) & \textbf{Params} & \textbf{Train Loss} & \textbf{Val PPL} & $\boldsymbol{\kl}$ \textbf{(nats)} & \textbf{Top-1} \\
\midrule
64 & 268,352 & 116.74 & 30.06 & 1.5824 & 60.3\% \\
128 & 530,560 & 98.57 & 19.56 & 1.2213 & 67.2\% \\
256 & 1,054,976 & 82.83 & 14.15 & 1.0029 & 69.3\% \\
512 & 2,103,808 & 68.90 & 8.91 & 0.7550 & 77.0\% \\
1024 & 4,201,472 & 52.56 & 6.96 & 0.7057 & 77.8\% \\
\textbf{2048} & \textbf{8,396,800} & \textbf{58.24} & \textbf{5.74} & \textbf{0.5783} & \textbf{81.7\%} \\
\bottomrule
\end{tabular}
\end{table}

\begin{figure}[t]
    \centering
    \includegraphics[width=0.75\textwidth]{figures/fig4_adapter_saturation.pdf}
    \caption{\textbf{Adapter capacity saturation.} Validation perplexity (purple, left axis) and KL divergence (red, right axis) vs.\ bottleneck rank $r$ at exit depth $L{=}18$ (50\%). The PPL curve exhibits asymptotic convergence toward the full-depth baseline (PPL$^*=5.76$, dashed green). While an $r{=}2048$ adapter adds only 17.2~MFLOPs per token ($< 0.6\%$ of the $\sim$3.0~GFLOPs saved by skipping 18 layers), Top-1 token agreement plateaus at 81.7\%, exposing an intrinsic expressivity and fidelity ceiling.}
    \label{fig:adapter}
\end{figure}

\paragraph{Structural interpretation and honest FLOP accounting.}
The capacity sweep reveals an asymptotic convergence pattern (Figure~\ref{fig:adapter}): validation PPL decreases from 30.06 ($r{=}64$) to 5.74 ($r{=}2048$). It is crucial to distinguish \emph{training cost} from \emph{inference operational FLOPs}:
\begin{itemize}[nosep,leftmargin=*]
    \item \textbf{Inference FLOPs are negligible:} For Qwen2.5-3B ($d_{\text{model}}=2048$), skipping 18 layers saves $\sim$3.0~GFLOPs per token. In contrast, an $r{=}512$ adapter incurs only $2 \times (2048 \times 512 + 512 \times 2048) \approx 4.2$~MFLOPs ($0.14\%$ of saved compute), and even an $r{=}2048$ adapter costs only $\approx 17.2$~MFLOPs ($0.57\%$ of savings). At inference, auxiliary adapter heads do \emph{not} eliminate compute savings.
    \item \textbf{The true barrier is the expressivity ceiling:} Despite negligible FLOPs, even an 8.4M-parameter adapter achieves only 81.7\% Top-1 token agreement and $\kl = 0.5783$~nats. Roughly 1 in 5 tokens remains mispredicted because intermediate states occupy an intrinsically distinct semantic manifold.
\end{itemize}

\paragraph{The core architectural asymmetry.}
This establishes the foundational insight of our investigation:

\begin{center}
\begin{tabular}{@{}lccl@{}}
\toprule
\textbf{Error Source} & \textbf{Effective Remedy} & \textbf{Learned Params} & \textbf{Inference FLOP Overhead} \\
\midrule
$\Eattn$ (KV cache continuity) & Native Projection & \textbf{0} & \textbf{0 FLOPs} (reuses skipped weights) \\
$\Ereadout$ (readout calibration) & Tuned Lens / Bottleneck & $0.26\text{M} - 2.10\text{M}$ & $< 4.2\text{ MFLOPs}$ ($< 0.15\%$ of savings) \\
\bottomrule
\end{tabular}
\end{center}

KV-cache continuity is solved exactly with zero learning and zero parameter overhead via Native Projection. For draft generation, pairing Native Projection with a lightweight calibration head provides a highly efficient production pipeline.

% ================================================================
% 7. MECHANISTIC EXPLANATION
% ================================================================
\section{Mechanistic Explanation: Key/Value Asymmetry}
\label{sec:mechanism}

The asymmetric difficulty of $\Eattn$ vs.\ $\Ereadout$ demands a mechanistic explanation. We hypothesize that it arises from a fundamental asymmetry in how Key and Value representations evolve across depth.

\subsection{Layer-wise Cosine Similarity Analysis}

For each exit depth $\dt$ and each subsequent layer $l > \dt$, we compute the cosine similarity between the oracle Key/Value states ($K_l^*, V_l^*$, computed from the true hidden state $h_l$) and the projected approximations ($\Khat_l, \Vhat_l$, computed from $h_{\dt}$):
\begin{equation}
    \cos(K_l^*, \Khat_l) = \frac{K_l^* \cdot \Khat_l}{\|K_l^*\| \|\Khat_l\|}, \quad
    \cos(V_l^*, \Vhat_l) = \frac{V_l^* \cdot \Vhat_l}{\|V_l^*\| \|\Vhat_l\|}.
\end{equation}

\begin{figure}[t]
    \centering
    \includegraphics[width=0.95\textwidth]{figures/fig3_key_vs_value_drift.pdf}
    \caption{\textbf{Key-stable / Value-evolving asymmetry.} Cosine similarity between oracle and projected Key (blue) and Value (red) states across layers beyond exit depth on Qwen2.5-3B, shown for three exit depths: $L{=}18$ (50\%), $L{=}24$ (67\%), and $L{=}30$ (83\%). Key similarity remains consistently high ($\ge 0.75$) across all skipped layers at every exit depth; Value similarity collapses rapidly. This explains why attention routing survives truncation while semantic readout does not.}
    \label{fig:kv_drift}
\end{figure}

\subsection{Empirical Findings}

Figure~\ref{fig:kv_drift} presents the layer-wise cosine similarity for exit depth $L{=}18$ (50\% of Qwen2.5-3B). The results reveal a striking asymmetry:

\begin{itemize}[nosep,leftmargin=*]
    \item \textbf{Key stability:} $\cos(K_l^*, \Khat_l)$ remains in the range $[0.75, 0.98]$ across all 18 skipped layers. The Key representations computed from the exit hidden state closely approximate the oracle Keys that would have been computed from the true (much deeper) hidden states.
    
    \item \textbf{Value instability:} $\cos(V_l^*, \Vhat_l)$ drops from $1.0$ at the exit layer to approximately $0.30$ within 2--3 layers, and remains in the range $[0.17, 0.30]$ for all subsequent layers. The projected Values are nearly orthogonal to the oracle Values.
\end{itemize}

This pattern is consistent across all exit depths and across both Qwen2.5-3B and Llama-3.2-1B (see Appendix), confirming its generality.

\subsection{Functional Interpretation}

The Key/Value asymmetry has a natural functional interpretation grounded in the Transformer's computation:

\begin{itemize}[nosep,leftmargin=*]
    \item \textbf{Keys govern attention routing:} $K_l$ determines \emph{which} tokens attend to \emph{which} other tokens. Attention routing patterns---particularly those related to syntactic structure, positional relationships, and in-context learning---are known to crystallize in early-to-mid layers~\cite{olsson2022context}. Once the routing topology stabilizes, subsequent layers' Key projections applied to a slightly ``stale'' hidden state still produce geometrically similar attention patterns.
    
    \item \textbf{Values carry semantic content:} $V_l$ carries the lexical, factual, and semantic information that is continuously refined by the feed-forward (MLP) sub-layers of the Transformer. As established by Geva et al.~\cite{geva2022transformer}, FFN layers function as key--value memories that progressively update the residual stream with factual knowledge. The residual stream evolves substantially through upper-layer MLPs:
    \begin{equation}
        h_l = h_{\dt} + \sum_{j=\dt+1}^{l} \left(\Delta h_{\text{attn}}^{(j)} + \Delta h_{\text{mlp}}^{(j)}\right),
    \end{equation}
    and Value projections $W_V^{(l)} \LN_1^{(l)}(\cdot)$ faithfully track this evolution. When applied to the stale $h_{\dt}$ instead of the true $h_l$, the resulting Values are qualitatively wrong.
\end{itemize}

\paragraph{Why Native Projection succeeds despite Value inaccuracy.}
Native Projection succeeds because multi-head attention is a \emph{weighted average} of Values, with weights determined by \emph{Keys}. If the Key-derived attention weights are accurate (high Key cosine similarity), the attention output is a correctly-weighted average---even if the individual Values are approximate. The attention mechanism is robust to Value perturbation when the routing structure is preserved, precisely because it was trained to be a soft retrieval mechanism over a continuously evolving value store.

% ================================================================
% 8. APPLICATION: SELF-SPECULATIVE DECODING
% ================================================================
\section{Application: Self-Speculative Decoding}
\label{sec:speculative}

To demonstrate the practical implications of KV-cache continuity, we implement a self-speculative decoding engine that uses the model's own first $\dt$ layers as a draft generator, with Native Projection maintaining cache continuity during speculation.

\subsection{Theoretical Framework: Naive vs.\ Partial-Layer Verification}

In self-speculative decoding, let $c = \dt / L$ denote the relative draft depth. We distinguish two distinct verification regimes:

\paragraph{Regime 1: Naive Full Verification.}
When the target verifier re-executes all $L$ layers from scratch, the expected speedup at horizon $\gamma{=}1$ is:
\begin{equation}
    \mathcal{S}_{\text{naive}} = \frac{1 + \alpha}{c + 1}, \quad \text{requiring } \alpha > c \text{ for } \mathcal{S}_{\text{naive}} > 1.
    \label{eq:speedup_naive}
\end{equation}
At $d_t = 24$ on Qwen2.5-3B ($c = 24/36 \approx 0.667$), break-even requires $\alpha > 66.7\%$. Because zero-shot intermediate readout achieves only $\alpha = 24.70\%$, naive verification yields $\mathcal{S}_{\text{naive}} = 0.748\times < 1.0\times$ (a slowdown).

\paragraph{Regime 2: Partial-Layer Verification via Native Projection.}
In self-speculative decoding, layers $1 \dots \dt$ have \emph{already} been executed during drafting, and their prefix states $h_{\dt}$ are stored in memory. By populating upper-layer KV caches ($l > \dt$) via Native Projection (Eq.~\ref{eq:kproj}--\ref{eq:vproj}), downstream attention continuity is preserved. The verification pass can therefore resume directly from $h_{\dt}$, evaluating \emph{only} the remaining $(L - \dt)$ layers. The effective verification cost drops from $1$ to $(1 - c)$:
\begin{equation}
    \mathcal{S}_{\text{partial}} = \frac{1 + \alpha}{c + (1 - c)} = \frac{1 + \alpha}{1} = \mathbf{1 + \alpha}.
    \label{eq:speedup_partial}
\end{equation}
Because $c + (1 - c) \equiv 1$, $\mathcal{S}_{\text{partial}} > 1$ for \emph{any} non-zero acceptance rate $\alpha > 0$. At $\alpha = 24.70\%$, partial-layer verification theoretically provides $\mathcal{S}_{\text{partial}} = \mathbf{1.247\times}$ net speedup, and lifting $\alpha \to 35\%$ via a lightweight calibrated readout head scales throughput to $\mathbf{1.35\times}$. Native Projection provides the mathematical bridge that makes partial-layer verification possible without model retraining.

\subsection{Large-Scale Benchmark}

We evaluate on $N{=}60$ prompts across 4 domains (Mathematics, Code, Science, Language; 15 prompts each) with $\dt{=}24$ and $\gamma{=}1$ on Qwen2.5-3B. Table~\ref{tab:speculative} summarizes the results.

\begin{table}[t]
\centering
\caption{\textbf{Multi-domain self-speculative decoding benchmark} ($N{=}60$, Qwen2.5-3B, $\dt{=}24$, $\gamma{=}1$). 95\% confidence intervals computed as $\mu \pm 1.96\sigma/\sqrt{N}$. Peak VRAM: 5.86 GB.}
\label{tab:speculative}
\small
\begin{tabular}{@{}lcccc@{}}
\toprule
\textbf{Domain} & $N$ & \textbf{Speedup} ($\mu \pm \text{CI}_{95}$) & $\boldsymbol{\alpha}$ ($\mu \pm \text{CI}_{95}$) & \textbf{Exact Match} \\
\midrule
Mathematics & 15 & $0.778 \pm 0.050\times$ & $28.56\% \pm 6.67\%$ & 93.3\% (14/15) \\
Code Generation & 15 & $0.669 \pm 0.028\times$ & $11.15\% \pm 4.86\%$ & 100.0\% (15/15) \\
Factual Knowledge & 15 & $0.785 \pm 0.026\times$ & $31.33\% \pm 4.22\%$ & 100.0\% (15/15) \\
Language \& Reasoning & 15 & $0.760 \pm 0.029\times$ & $27.75\% \pm 5.84\%$ & 93.3\% (14/15) \\
\midrule
\textbf{All Domains} & \textbf{60} & $\mathbf{0.748 \pm 0.021\times}$ & $\mathbf{24.70\% \pm 3.35\%}$ & \textbf{96.67\% (58/60)} \\
\bottomrule
\end{tabular}
\end{table}

\paragraph{Lossless token equivalence.}
The $96.67\%$ exact match rate (58/60 prompts producing identical token sequences to standard autoregressive decoding) validates that Native Projection preserves generation semantics. The measured speedup of $0.748\times$ exactly matches the theoretical prediction of $\mathcal{S} = (1+\alpha)/(c+1) = (1+0.247)/(0.667+1) = 0.748\times$ for $\gamma{=}1$. The sub-unity speedup is a direct consequence of the high cost ratio $c = \dt/L = 0.667$: at $\gamma{=}1$ the denominator $(c + 1) = 1.667$ exceeds the expected tokens per round $(1+\alpha) = 1.247$, and increasing $\gamma$ does not help---at $\gamma{=}3$ the formula yields $\mathcal{S} \approx 0.44\times$, since the wasted draft compute ($c\gamma$) grows faster than the accepted token count ($\approx 1.32$) at this acceptance rate. Practical super-unity speedup from self-speculative decoding therefore requires either a \emph{shallower} draft depth (lower $c$) at the cost of reduced acceptance rate, or amortized verification via partial-layer caching~\cite{elhoushi2024layerskip} that reduces verification cost from $1$ to $(1{-}c)$. The primary contribution of Native Projection in this context is not raw throughput but \emph{cache continuity}: ensuring that speculative tokens do not poison the KV cache for subsequent generation.

\begin{figure}[t]
    \centering
    \includegraphics[width=0.95\textwidth]{figures/fig6_speculative_benchmark.pdf}
    \caption{\textbf{Self-speculative decoding domain breakdown} ($N{=}60$, Qwen2.5-3B, $\dt{=}24$, $\gamma{=}1$). (a)~Measured speedup per domain with 95\% confidence intervals. Aggregate speedup of $0.751\times$ matches the theoretical prediction of $0.748\times$. (b)~Per-domain acceptance rates with exact-match counts shown above each bar. Code generation has the lowest acceptance rate (11.5\%) but achieves 100\% exact match (15/15).}
    \label{fig:speculative}
\end{figure}

\paragraph{Memory efficiency.}
Peak VRAM consumption is 5.86~GB, using only 39\% of a 15~GB consumer GPU budget. Because there is only one model instance (the draft and verifier share all weights), memory overhead beyond model weights is limited to the static KV cache (72~MB for 2{,}048 tokens across 36 layers with GQA). Native Projection itself requires zero additional parameters.

% ================================================================
% 9. DISCUSSION
% ================================================================
\section{Discussion}
\label{sec:discussion}

\paragraph{Implications for depth-adaptive inference.}
Our results reframe the challenge of depth-adaptive Transformer inference. The common assumption that ``early exit quality degrades because the model hasn't finished processing'' is imprecise: it is specifically the \emph{readout} that fails, not the \emph{attention memory}. This distinction has immediate practical implications:

\begin{enumerate}[nosep,leftmargin=*]
    \item \textbf{For speculative decoding:} KV-cache continuity via Native Projection enables zero-recomputation verification, eliminating the primary overhead of self-speculative approaches like LayerSkip~\cite{elhoushi2024layerskip}.
    
    \item \textbf{For early exit systems:} The negative result on standalone readout implies that lightweight exit heads (as in CALM~\cite{schuster2022confident} or DeeBERT~\cite{xin2020deebert}) cannot match full-depth quality for autoregressive generation without substantial auxiliary capacity, calling into question the viability of hard early exit for generation tasks.
    
    \item \textbf{For model compression:} The Key-stable/Value-evolving asymmetry suggests that depth-wise compression strategies should preserve Key projection fidelity above all else, as attention routing is the more fragile component to approximate poorly.
\end{enumerate}

\paragraph{The 61\% boundary.}
The observation that Native Projection crosses the 1-bit boundary at $\sim$61\% depth on Qwen2.5-3B and $\sim$50\% on Llama-3.2-1B suggests a general property of pre-trained Transformer residual streams: the hidden state stabilizes sufficiently for accurate Key projection in the second half of the network. We hypothesize this relates to the ``representation convergence'' phenomenon observed in deep networks~\cite{raghu2017svcca}, where upper layers perform increasingly fine-grained refinements rather than qualitative representation changes.

% ================================================================
% 10. LIMITATIONS
% ================================================================
\section{Limitations}
\label{sec:limitations}

\begin{enumerate}[nosep,leftmargin=*]
    \item \textbf{Model scale:} Our experiments use 1B and 3B parameter models. Validation on 7B+ and MoE architectures would strengthen the generality claims.
    
    \item \textbf{Evaluation scope:} All quantitative results use WikiText-2 under teacher forcing. While the speculative decoding benchmark spans 4 domains with free generation, broader evaluation on downstream NLU/NLG benchmarks would be valuable.
    
    \item \textbf{Adapter training budget:} The capacity ablation trains each adapter for 3{,}000 steps. While the asymptotic convergence pattern (Figure~\ref{fig:adapter}) strongly suggests a structural rather than optimization limitation, longer training schedules could further tighten this bound.
    
    \item \textbf{Speculation horizon:} The live benchmark uses $\gamma{=}1$ with $c{=}0.667$, yielding sub-unity speedup. Practical super-unity speedup requires either a shallower draft depth (lower $c$), higher acceptance rates, or amortized verification via partial-layer caching.
\end{enumerate}

% ================================================================
% 11. CONCLUSION
% ================================================================
\section{Conclusion}

We have introduced a causal disentanglement protocol that cleanly separates two failure modes of depth-truncated Transformer inference: attention memory poisoning and unembedding manifold mismatch. Through this protocol, we established that:

\begin{enumerate}[nosep,leftmargin=*]
    \item \textbf{KV-cache continuity is solved without training}: Native KV State Projection preserves multi-head attention quality across skipped layers for exit depths $\ge 61\%$, using only the model's existing frozen weights.
    
    \item \textbf{Standalone early exit readout is structurally non-viable}: Matching full-depth quality from intermediate representations requires adapter capacity ($r{=}2048$, 8.4M parameters) that eliminates the compute savings of early exit.
    
    \item \textbf{A Key-stable / Value-evolving asymmetry} explains both results mechanistically: attention routing geometry crystallizes early while semantic content continues evolving through upper-layer MLPs.
\end{enumerate}

These findings inform the design of future depth-adaptive inference systems and explain why the field has moved from self-speculative approaches toward dedicated draft models. We release our disentanglement harness and evaluation code to enable the community to apply this diagnostic framework to new architectures and adaptive inference strategies.

% ================================================================
% REFERENCES
% ================================================================
\bibliographystyle{plainnat}

\begin{thebibliography}{25}

\bibitem[Belrose et al.(2023)]{belrose2023eliciting}
N.~Belrose, Z.~Furman, L.~Smith, D.~Halawi, I.~Ostrovsky, L.~McKinney, S.~Biderman, and J.~Steinhardt.
\newblock Eliciting latent predictions from transformers with the tuned lens.
\newblock \emph{arXiv preprint arXiv:2303.08112}, 2023.

\bibitem[Cai et al.(2024)]{cai2024medusa}
T.~Cai, Y.~Li, Z.~Geng, H.~Peng, J.~D. Lee, D.~Chen, and T.~Dao.
\newblock Medusa: Simple {LLM} inference acceleration framework with multiple decoding heads.
\newblock In \emph{ICML}, 2024.

\bibitem[Chen et al.(2023)]{chen2023accelerating}
C.~Chen, S.~Borgeaud, G.~Irving, J.-B. Lespiau, L.~Sifre, and J.~Jumper.
\newblock Accelerating large language model decoding with speculative sampling.
\newblock \emph{arXiv preprint arXiv:2302.01318}, 2023.

\bibitem[Del~Corro et al.(2023)]{del2023skipdecode}
L.~Del~Corro, A.~Del~Giorno, S.~Alon, and A.~Sabharwal.
\newblock {SkipDecode}: Autoregressive skip decoding with batched verification.
\newblock \emph{arXiv preprint arXiv:2307.02628}, 2023.

\bibitem[Elbayad et al.(2020)]{elbayad2020depth}
M.~Elbayad, J.~Gu, E.~Grave, and M.~Auli.
\newblock Depth-adaptive transformer.
\newblock In \emph{ICLR}, 2020.

\bibitem[Elhoushi et al.(2024)]{elhoushi2024layerskip}
M.~Elhoushi, A.~Shrivastava, D.~Liskovich, B.~Hosmer, B.~Wasti, L.~Lai, A.~Ber, C.~Kim, A.~Bhatia, S.~Mahmoud, et~al.
\newblock {LayerSkip}: Enabling early exit inference and self-speculative decoding.
\newblock In \emph{ACL}, 2024.

\bibitem[Geva et al.(2022)]{geva2022transformer}
M.~Geva, A.~Caciularu, K.~Wang, and Y.~Goldberg.
\newblock Transformer feed-forward layers build predictions by promoting concepts in the vocabulary space.
\newblock In \emph{EMNLP}, 2022.

\bibitem[Leviathan et al.(2023)]{leviathan2023fast}
Y.~Leviathan, M.~Kalman, and Y.~Matias.
\newblock Fast inference from transformers via speculative decoding.
\newblock In \emph{ICML}, 2023.

\bibitem[Li et al.(2024)]{li2024eagle}
Y.~Li, F.~Wei, C.~Zhang, and H.~Zhang.
\newblock {EAGLE}: Speculative sampling requires rethinking feature uncertainty.
\newblock In \emph{ICML}, 2024.

\bibitem[Merity et al.(2017)]{merity2017pointer}
S.~Merity, C.~Xiong, J.~Bradbury, and R.~Socher.
\newblock Pointer sentinel mixture models.
\newblock In \emph{ICLR}, 2017.

\bibitem[Nostalgebraist(2020)]{nostalgebraist2020logitlens}
Nostalgebraist.
\newblock Interpreting {GPT}: the logit lens.
\newblock \emph{LessWrong Blog Post}, 2020.

\bibitem[Olsson et al.(2022)]{olsson2022context}
C.~Olsson, N.~Elhage, N.~Nanda, N.~Joseph, N.~DasSarma, T.~Henighan, B.~Mann, A.~Askell, Y.~Bai, A.~Chen, et~al.
\newblock In-context learning and induction heads.
\newblock \emph{arXiv preprint arXiv:2209.11895}, 2022.

\bibitem[Raghu et al.(2017)]{raghu2017svcca}
M.~Raghu, J.~Gilmer, J.~Yosinski, and J.~Sohl-Dickstein.
\newblock {SVCCA}: Singular vector canonical correlation analysis for deep learning dynamics and interpretability.
\newblock In \emph{NeurIPS}, 2017.

\bibitem[Raposo et al.(2024)]{raposo2024mixture}
D.~Raposo, S.~Ritter, B.~Richards, T.~Lillicrap, P.~Conway~Humphreys, and A.~Santoro.
\newblock Mixture-of-depths: Dynamically allocating compute in transformer-based language models.
\newblock \emph{arXiv preprint arXiv:2404.02258}, 2024.

\bibitem[Schuster et al.(2022)]{schuster2022confident}
T.~Schuster, A.~Fisch, J.~Jaakkola, and R.~Barzilay.
\newblock Confident adaptive language modeling.
\newblock In \emph{NeurIPS}, 2022.

\bibitem[Schwartz et al.(2020)]{schwartz2020right}
R.~Schwartz, G.~Stanovsky, S.~Swayamdipta, J.~Dodge, and N.~A. Smith.
\newblock The right tool for the job: Matching model and instance complexities.
\newblock In \emph{ACL}, 2020.

\bibitem[Xiao et al.(2024)]{xiao2024efficient}
G.~Xiao, Y.~Tian, B.~Chen, S.~Han, and M.~Lewis.
\newblock Efficient streaming language models with attention sinks.
\newblock In \emph{ICLR}, 2024.

\bibitem[Xin et al.(2020)]{xin2020deebert}
J.~Xin, R.~Tang, J.~Lee, Y.~Yu, and J.~Lin.
\newblock {DeeBERT}: Dynamic early exiting for accelerating {BERT} inference.
\newblock In \emph{ACL}, 2020.

\bibitem[Zhang et al.(2023)]{zhang2023h2o}
Z.~Zhang, Y.~Sheng, T.~Zhou, T.~Chen, L.~Zheng, R.~Cai, Z.~Song, Y.~Tian, C.~Ré, C.~Barrett, Z.~Wang, and B.~Chen.
\newblock {H$_2$O}: Heavy-hitter oracle for efficient generative inference of large language models.
\newblock In \emph{NeurIPS}, 2023.

\bibitem[Zhang et al.(2024)]{zhang2024draft}
J.~Zhang, J.~Wang, H.~Li, L.~Shou, K.~Chen, G.~Chen, and S.~Mehrotra.
\newblock Draft \& verify: Lossless large language model acceleration via self-speculative decoding.
\newblock In \emph{ACL}, 2024.

\end{thebibliography}

% ================================================================
% APPENDIX
% ================================================================
\newpage
\appendix

\section{Complete Qwen2.5-3B Depth Sweep Data}
\label{app:full_sweep}

Table~\ref{tab:full_sweep} presents the complete 11-point depth sweep for Qwen2.5-3B across all conditions. Data sourced from canonical evaluation on WikiText-2 ($T{=}64$, teacher-forced).

\begin{table}[h]
\centering
\caption{\textbf{Complete depth sweep} on Qwen2.5-3B ($\text{PPL}^* = 5.76$, 36 layers). Native Projection vs.\ Direct Copy across 11 exit depths. The $\ln(2)$ crossing occurs between $L{=}18$ and $L{=}22$ for Native Projection.}
\label{tab:full_sweep}
\small
\begin{tabular}{@{}cccccc@{}}
\toprule
\textbf{Exit} & \textbf{Rel.} & \multicolumn{2}{c}{\textbf{Native Projection}} & \multicolumn{2}{c}{\textbf{Direct Copy}} \\
\cmidrule(lr){3-4} \cmidrule(lr){5-6}
$L$ & \textbf{Depth} & PPL & $\kl$ (nats) & PPL & $\kl$ (nats) \\
\midrule
6 & 16.7\% & 43.22 & 2.1371 & 245,874 & 10.5297 \\
10 & 27.8\% & 28.72 & 1.7265 & 154,360 & 10.1770 \\
14 & 38.9\% & 20.42 & 1.4110 & 123,257 & 9.6513 \\
18 & 50.0\% & 13.33 & 0.9604 & 3,088 & 6.3036 \\
22 & 61.1\% & 9.90 & 0.6151 & 4,276 & 6.3835 \\
24 & 66.7\% & 8.04 & 0.4365 & 3,450 & 6.0820 \\
26 & 72.2\% & 7.41 & 0.2929 & 6,133 & 6.7878 \\
28 & 77.8\% & 6.86 & 0.1692 & 4,349 & 5.9263 \\
30 & 83.3\% & 6.45 & 0.1049 & 10,274 & 5.8089 \\
32 & 88.9\% & 6.16 & 0.0630 & 57.40 & 2.4687 \\
34 & 94.4\% & 5.94 & 0.0194 & 9.58 & 0.5946 \\
\bottomrule
\end{tabular}
\end{table}

\section{Llama-3.2-1B Cosine Similarity Data}
\label{app:llama_cosine}

Table~\ref{tab:llama_cosine} presents representative layer-wise cosine similarity values for Llama-3.2-1B at exit depth $L{=}8$ (50\%), confirming the Key-stable/Value-evolving asymmetry generalizes across architectures.

\begin{table}[h]
\centering
\caption{\textbf{Key vs.\ Value cosine similarity} on Llama-3.2-1B at exit depth $L{=}8$ (50\%, 16 layers). Key similarity remains $\ge 0.63$ while Value similarity drops to $\le 0.59$ by 2 layers beyond exit.}
\label{tab:llama_cosine}
\small
\begin{tabular}{@{}cccccc@{}}
\toprule
\textbf{Layer} & $\cos(K^*, \Khat)$ & $\cos(V^*, \Vhat)$ \\
\midrule
9 (exit+1) & 0.9668 & 0.7713 \\
10 (exit+2) & 0.9195 & 0.7358 \\
11 (exit+3) & 0.8400 & 0.5696 \\
12 (exit+4) & 0.8087 & 0.5549 \\
13 (exit+5) & 0.7654 & 0.5267 \\
14 (exit+6) & 0.7331 & 0.5902 \\
15 (exit+7) & 0.6330 & 0.5375 \\
\bottomrule
\end{tabular}
\end{table}

\section{Self-Speculative Decoding Engine Details}
\label{app:engine}

\paragraph{Architecture.}
The self-speculative decoding engine operates with a single model instance. During the draft phase, only layers $1$ through $\dt$ are executed, with Native Projection populating cache slots for layers $\dt{+}1$ through $L$. During verification, only the upper layers $\dt{+}1$ through $L$ execute on the draft tokens as a single batched forward pass, reusing the pristine cache entries for layers $1$ through $\dt$ written during drafting. This achieves \textbf{zero cache recomputation}---a key advantage over LayerSkip~\cite{elhoushi2024layerskip}, which requires full upper-layer recomputation during verification.

\paragraph{Memory budget.}
For Qwen2.5-3B with GQA ($n_{\text{kv}} = 2$, $d_{\text{head}} = 128$):
\begin{align}
    M_{\text{weights}} &= 3.09 \times 10^9 \times 2\;\text{bytes} = 6.18\;\text{GB}, \\
    M_{\text{KV cache}} &= 36 \times 2 \times 2 \times 2048 \times 128 \times 2\;\text{bytes} = 72.35\;\text{MB}, \\
    M_{\text{peak}} &\le 5.86\;\text{GB} \quad \text{(measured)}.
\end{align}
Native Projection adds no learned parameters; its compute cost is limited to the projection FLOPs detailed in Section~\ref{sec:method}.

\section{Reproducibility}
\label{app:reproducibility}

All experiments use:
\begin{itemize}[nosep,leftmargin=*]
    \item \textbf{Evaluation dataset:} WikiText-2 (raw, validation split)
    \item \textbf{Sequence length:} $T = 64$ tokens (teacher-forced disentanglement); up to 30 tokens (free generation benchmark)
    \item \textbf{Random seed:} 42
    \item \textbf{Precision:} \texttt{bfloat16}
    \item \textbf{Hardware:} Single NVIDIA GPU ($\le$ 15 GB VRAM)
\end{itemize}

All canonical results are stored in CSV format under \texttt{results/canonical\_sweeps/} with full metadata. The disentanglement harness is provided in \texttt{experiments/isolate\_kv\_ablation.py}.

\end{document}
